\section{A Type-Theoretic Approach}\label{sec:type-theory}

Consider the addition of recursive descriptors to the Pollux formalization. With
a descriptor rule like this one
\[
  \infer[D-Update]{ \dom{d} \subseteq \dom{d'} \\ \forall\ k \in
    \dom{d} \cap \dom{d'}.\ d[k] \propto d'[k] }{ d \ll d' }
\]
the descriptor relation would become co-inductive since an update to this
descriptor would would require showing $d_1[k] \propto d_2[k]$ if $k$ is the field
were we see the recursive type and the field relation would include:
\[ \infer[F-Msg]{ d_1[k] \ll d_2[k] }{ \langle k, d_1[k] \rangle \propto \langle k, d_2[k] \rangle } \]
So a derivation that $d_1[k] \ll d_2[k]$ would need to include a derivation of
itself! The natural solution here is to make $\ll$ co-inductive. 

However, if we more closely model the descriptors as proper recursive types, it
might be possible to avoid that. Consider this Protobuf descriptor:
\begin{minted}{proto}
message Stream {
  int hd = 1;
  Stream tl = 2;
}
\end{minted}
Note first that while it's called \texttt{Stream}, since all Protobuf fields can
be omitted and message fields always have an explicit presence indicator, we
actually have a \texttt{List} definition rather than a true co-inductive stream.

From the prospective of algebraic type theory, this type is $F\ X = \mathtt{int}
+ \mathtt{int} \times F\ X$. If we were to add a \texttt{bool} field to this
descriptor, the algebraic type would update to $G\ X = \mathtt{int} \times
\mathtt{bool} + \mathtt{int} \times \mathtt{bool} \times G\ X$. 

Then the descriptor relation I'm proposing would be a small proof system where
we want to show that two descriptors are compatible if $\varnothing \vdash d \ll d'$, but we could
\emph{define} what it means for two descriptor to be compatible as 
\begin{mathpar}
  d_1 := \mu\ X.\ F\ X 

  d_2 := \mu\ X.\ G\ X 

  d_1 \ll d_2 := \exists\ f: A \rightarrow B.\ \forall\ m \checkmark^{d_1}.\ \exists\ bs.\ \mathrm{Encodes}_{d_1}(m, bs) = bs \wedge
  \mathrm{parse}_{d_2}(bs) = f(m)
\end{mathpar}

In English, then the definition of two descriptors being compatible is that
there exists a transformation such that for all valid messages under descriptor
$d_1$, we can serialize that message into some encoding that can be parsed to
produce the image of $m$ under the transformation. Rather than trying to prove
$\varnothing \vdash d_1 \ll d_2$, which would require co-induction, we show that 
\[ \mu\ F \ll \mu\ G \vdash F (\mu\ F) \ll G (\mu\ G) \]
which can be done with normal induction. This should be equivalent to the
original statement using the relation-based definition, when viewed through the
right lens of folding and unfolding recursive types.

If the descriptor relation is a static typing judgment, then the transformation
is the semantic type judgment. This would allow us to avoid co-induction, 

However, we even if we present this framework first, we probably still want to
have a relation on the descriptors. The above definition doesn't provide any
insight into how to write the transformation, while we've observed that having
the relation provides a structure that the transformation mirrors. Additionally,
I personally think that the relation on descriptors provides a nice way to check
compatibility without having to construct the transformation directly.

We can take this view even farther though, and make the argument that Pollux
become a type system which can characterize functions over messages has a
compatible transformation.

\[ \vDash f : d_1 \preceq d_2 \rightarrow \forall\ m \checkmark^{d_1}.\ \exists\ bs.\ \mathrm{Encodes}_{d_1}(m, bs) \wedge
  \mathrm{parse}_{d_2}(bs) = f(m) \]

Then we end up with a syntactic typing judgment on the descriptors themselves,
related to what we can convert between, and a semantic typing judgment on the
transformation between the descriptors. Pulling in the recursive type view of
the descriptors, we'd end up with a typing rule something like this (likely
\emph{not} actually correct):

\[ \infer{ X : \mu\ F \preceq \mu\ G \vdash f : F(\mu\ F) \preceq G(\mu\ G) }{\vDash \mu\ X.f : \mu\ F \preceq \mu\ G} \]

This is basically a logical relation in disguise. We want to prove something
about the transformation functions (which is an infinite family of functions),
and to do so we build the descriptor relation as a guide.

% Local Variables:
% citar-bibliography: ("../pollux.bib")
% TeX-master: "../pollux.tex"
% TeX-command-default: "Make"
% jinx-local-words: "protobuf"
% End:
